% DreamMachine -- single-weakness study. DRAFT v0.1 (2026-11-10).
% Every number comes from the repository: docs/RUNNING.md sections 10-15 and 17, and
% runs/pipelines/{80b37628b193,409305f3e2f2}/results.json. Items marked [TODO] need the team.
\documentclass[11pt]{article}
\usepackage[margin=1in]{geometry}
\usepackage[T1]{fontenc}
\usepackage{lmodern}
\usepackage{amsmath,amssymb}
\usepackage{booktabs}
\usepackage{multirow}
\usepackage{graphicx}
\usepackage{xcolor}
\usepackage[hidelinks]{hyperref}
\usepackage[numbers,sort&compress]{natbib}

\newcommand{\todo}[1]{\textcolor{red}{[TODO: #1]}}
\newcommand{\eg}{e.g.,\ }
\newcommand{\ie}{i.e.,\ }

\title{Targeting or Difficulty? A Controlled Test of Weakness-Targeted\\
Synthetic Fine-Tuning for Small Language Models}
\author{\todo{author names and affiliations}}
\date{Draft, 10 November 2026}

\begin{document}
\maketitle

\begin{abstract}
Recent methods diagnose where a language model fails and then fine-tune it on data built around those failures.
Such data is, almost by definition, harder for the model than the data it is compared with, so an observed gain
may come from difficulty rather than from aiming at the right skill. We separate the two. A program generates
arithmetic word problems from computation graphs, so every problem's skill profile (number of steps, operation
counts, magnitude, distractors) is known exactly instead of being tagged by another model. A Linear Logistic Test
Model fitted to the model's answers on a fully crossed probe set ranks these factors by their effect on accuracy.
We then train four arms at an equal token budget: real problems only, untargeted generated problems, problems aimed
at the strongest weakness, and a control whose problems are matched one-to-one on predicted difficulty but are hard
for a different reason. On Qwen3-0.6B and Qwen3-1.7B, fine-tuned with QLoRA on a single 8\,GB laptop GPU (three
seeds each), both models show the same weakness (long multi-step problems), and aiming at it shrinks that factor's
estimated effect more than any other arm. Yet the targeted arm is no better than the matched control on GSM8K or
GSM-Symbolic (differences between $-0.012$ and $+0.015$, all intervals spanning zero), and on held-out generated
problems it is \emph{worse} ($-0.158$ for 0.6B, $-0.047$ for 1.7B). Every fine-tuned arm, including the one trained
only on real GSM8K problems, loses 14--24 points on GSM8K. A paired answer-level analysis traces this loss to the
terse style of the training solutions: the models stop restating the given facts and make more planning errors.
\end{abstract}

% ---------------------------------------------------------------------------------------------------------------
\section{Introduction}

Small open models (below about two billion parameters) can be adapted on consumer hardware, which makes them an
attractive target for cheap, focused improvement. A common recipe has two stages: first find out what the model
is bad at, then synthesise training examples that exercise exactly those skills
\citep{gupta2026craft,luo2026kite,liang2025sws,yu2025selferror}. Reported gains are usually measured against
randomly drawn or untargeted data.

That comparison has a built-in confound. Data selected because the model fails on it is harder for the model than
random data. Training on harder examples can help for reasons that have nothing to do with which skill they
exercise, for example because they carry more learning signal per token. Some methods already restrict training
problems to a band of pass rates \citep{liang2025sws}, which narrows the gap, but to our knowledge no study has
held difficulty fixed while varying only \emph{why} the problems are hard. Without that, a targeting method cannot
show that its targeting is what helps.

This paper runs that test for multi-step arithmetic word problems. Our design has three parts.
\begin{enumerate}
  \item \textbf{Exact skill profiles.} Every training and probe problem is produced from a computation graph by a
  program, so its features are computed rather than annotated by a language model.
  \item \textbf{A statistical diagnosis.} A Linear Logistic Test Model (LLTM) \citep{fischer1973lltm} fitted on a
  fully crossed probe set gives each factor's effect on the log-odds of a correct answer, with bootstrap intervals.
  \item \textbf{A difficulty-matched control.} For every targeted problem we pick an off-target problem with the
  same predicted difficulty under the fitted model, so the two arms differ only in which factor makes problems hard.
\end{enumerate}

\paragraph{Findings.} On two model sizes the answer is the same: targeting does what it says on the narrow measure
(the targeted factor's effect shrinks most), but it does not translate into better accuracy on real benchmarks,
and on held-out generated problems the off-target control does better. We also document a large accuracy cost
shared by every fine-tuned arm and give an answer-level account of where it comes from. All experiments run on one
RTX~4060 laptop GPU with 8\,GB of memory.

% ---------------------------------------------------------------------------------------------------------------
\section{Related Work}

\paragraph{Weakness-driven data synthesis.}
CRAFT clusters rubric criteria into a capability tree, scores a model per node and generates fine-tuning data for
the weakest nodes using a teacher model \citep{gupta2026craft}. \citet{luo2026kite} diagnose skills with a
cognitive diagnosis model over a skill matrix tagged by a language model, generate candidates and keep those within
an uncertainty band, in an iterative self-training loop. SwS identifies problems a model fails during reinforcement
learning and synthesises new problems from their concepts, keeping only those the model solves part of the time
\citep{liang2025sws}. Self-Error-Instruct groups a model's own error types and synthesises examples per type
\citep{yu2025selferror}. These methods differ in how they diagnose and generate, but they evaluate targeting mainly
against untargeted or random data. Our contribution is the difficulty-matched comparison, made possible by skill
profiles that are exact by construction.

\paragraph{Programmatic arithmetic data.}
\citet{gangwar2025arith} show that programmatically generated arithmetic data can improve the mathematical
reasoning of smaller models, either as an intermediate fine-tuning stage or mixed into instruction data. Our
generator is in this spirit, but we use it to diagnose and to control the data, not only to supply more of it.
GSM-Symbolic \citep{mirzadeh2024gsmsymbolic} builds templated variants of GSM8K \citep{cobbe2021gsm8k} to measure
robustness; we use it only as an evaluation set.

\paragraph{Item response models.}
The LLTM \citep{fischer1973lltm} explains item difficulty as a weighted sum of item properties; it belongs to the
family of explanatory item response models \citep{deboeck2004explanatory}. Because our probe items are generated
from a crossed design, the factor weights can be read as the effect of changing one property while the others are
held fixed.

\paragraph{Forgetting during fine-tuning.}
Narrow supervised fine-tuning can erode abilities a model had before \citep{kirkpatrick2017ewc}. Low-rank adapters
\citep{hu2021lora,dettmers2023qlora} reduce but do not remove this \citep{biderman2024lora}. Online reinforcement
learning appears to forget less than supervised learning on external text \citep{shenfeld2025razor}, and mixing in
the model's own generations counters forgetting largely by matching its style \citep{sun2026talking}. Rejection
sampling of a model's own solutions \citep{zelikman2022star,tong2024dart} is another way to keep training text close
to the model's distribution; \citet{tong2024dart} note that it under-represents hard problems unless more attempts
are spent on them. The observation that supervised fine-tuning mostly shapes the format and style of answers
\citep{zhou2023lima} is directly relevant to the accuracy cost we report in Section~\ref{sec:cost}.

% ---------------------------------------------------------------------------------------------------------------
\section{Method}

\subsection{Problems with exact skill profiles}
Each problem is a directed acyclic graph whose leaves are the numbers stated in the text and whose internal nodes
are the operations $+,-,\times,\div$, evaluated in exact rational arithmetic. A surface family (one of ten story
types, eight used for training and two held out) turns the graph into English sentences. From the graph we compute
a feature vector $q_i$ with the following entries: \texttt{steps} (number of operation nodes),
\texttt{log10\_max} ($\log_{10}$ of the largest intermediate value), \texttt{n\_mul} and \texttt{n\_div} (operation
counts), \texttt{n\_carry} (column carries and borrows), \texttt{n\_distractors} (numeric sentences not used in the
graph) and \texttt{has\_merge} (whether two separate chains are combined). These features are exact for every
generated problem.

\subsection{Diagnosis with a Linear Logistic Test Model}
The probe set crosses steps $\in\{2,\dots,7\}$, digits per number $\in\{1,\dots,4\}$, distractors
$\in\{0,1,2\}$ and three operation profiles (additive, mixed, multiplicative), giving 216 cells. We draw four
problems per cell (864 problems) and sample three answers per problem at temperature 0.7. With $s_i$ correct answers
out of $n_i$, the model is
\begin{equation}
  \operatorname{logit} P(\text{correct}_i) = \theta - \sum_k q_{ik}\,\eta_k ,
\end{equation}
where $\theta$ is the model's ability and $\eta_k$ is the cost in log-odds of one unit of factor $k$. We fit
$(\theta,\eta)$ by maximum likelihood with a small ridge penalty on $\eta$ and obtain 95\% intervals from 500
bootstrap resamples of the items \citep{efron1993bootstrap}. Because factors have different ranges, we rank them
by a range-normalised effect
\begin{equation}
  E_k = \eta_k\,\bigl(Q_k^{90} - Q_k^{10}\bigr),
\end{equation}
where $Q_k^{p}$ is the $p$-th percentile of factor $k$ over a pool of 5{,}000 problems drawn from a GSM8K-like
natural distribution. The target factor is the one with the largest $E_k$ among those whose interval excludes zero,
and the target threshold is the factor's natural median.

\subsection{Four training arms at an equal budget}
All arms receive 600{,}000 whitespace tokens of training text. Generated problems are mixed with real GSM8K
training problems at a ratio of 3:1 (synthetic:real); the targeted arm is also run at 9:1.
\begin{itemize}
  \item \textbf{real\_only}: GSM8K training problems only.
  \item \textbf{untargeted}: generated problems from the natural distribution.
  \item \textbf{targeted}: generated problems above the target threshold whose predicted probability of a correct
  answer, $\hat p_i$, lies in $[0.3,0.7]$, drawn with weight $\hat p_i(1-\hat p_i)$, the information a single
  Bernoulli answer carries about ability.
  \item \textbf{matched\_control}: generated problems at or below the target threshold. Each targeted problem is
  paired greedily, without replacement, with the unused off-target problem whose predicted logit is closest.
\end{itemize}
All candidates come from one generated pool of 40{,}000 problems. Match quality is reported as the mean absolute
logit difference and the two-sample Kolmogorov--Smirnov statistic between the arms' predicted accuracies. When an arm
would fall short of the token budget (the control's problems are shorter), the number of selected problems is sized
by the shorter side and the pool is enlarged deterministically.

\subsection{Error classification without annotation}
Every wrong answer is assigned one label, in this order: \textsc{format\_error} (no extractable answer),
\textsc{arithmetic\_slip} (a written equation is false), \textsc{unverifiable} (no equations to check),
\textsc{distractor\_use} (an operand equals an irrelevant number) and otherwise \textsc{plan\_error}. Reference
steps come from the graph for generated problems and from the calculator annotations for GSM8K. A human check of
these labels (Cohen's $\kappa$ on 150 stratified answers) is prepared but not yet done \todo{run the human
annotation and report $\kappa$}. GSM-Symbolic has no step annotations, so we report only its accuracy, not its error
types.

% ---------------------------------------------------------------------------------------------------------------
\section{Experimental Setup}

\paragraph{Models and prompt.} Qwen3-0.6B and Qwen3-1.7B \citep{qwen3}, thinking mode off, with the step-by-step
math instruction recommended on the Qwen3 model card and a boxed final answer. On a 200-problem development slice of
GSM8K-train this prompt was the most accurate of the three we tried for both models.

\paragraph{Training.} QLoRA \citep{dettmers2023qlora}: 4-bit NF4 base, rank 16, $\alpha=32$, dropout 0.05,
learning rate $2\times10^{-4}$, two epochs, effective batch 16, maximum length 512 tokens. Training solutions are
short: GSM8K reference solutions and, for generated problems, one line per computation step. Each arm is trained
from the base model with three seeds (0, 1, 2). One RTX~4060 laptop GPU (8\,GB); a training run took about
1.5\,h and an evaluation 15--26\,min.

\paragraph{Evaluation.} Greedy decoding, 512 new tokens, bf16 weights. After training, each model answers the first
300 GSM8K test problems, 200 GSM-Symbolic problems spread over its 100 templates, and 135 generated probe problems
from each of the training and held-out families. The untrained model answers a superset of these, so every
comparison is paired on the same problems. Differences are means over the three seeds with 95\% intervals from a
paired bootstrap over problems.

\paragraph{Protocol.} The configuration was fixed before the runs, each run started from a clean code version, and
the reported numbers come from those runs only. A separate development study (Section~\ref{sec:cost}) used GSM8K
training problems held out from training and is reported as exploratory.

% ---------------------------------------------------------------------------------------------------------------
\section{Results}

\subsection{Diagnosis: both models fail on long problems}
Table~\ref{tab:diag} lists the significant factors. For both models the largest range-normalised effect is the
number of steps, so the target is problems with more than four steps. The 0.6B model answers 36\% of the 864 probe
problems correctly and the 1.7B model 64\%. Accuracy falls with each added step: from 73\% at two steps to 14\% at
seven for 0.6B, and from 89\% to 37\% for 1.7B. Division is the second factor for both. Carries are not significant
for either model once the other factors are included, and the magnitude factor is significant only for 0.6B.
McFadden's pseudo-$R^2$ \citep{mcfadden1974} is 0.28 (0.6B) and 0.25 (1.7B).

\begin{table}[t]
\centering
\caption{LLTM diagnosis on 864 crossed probe problems $\times$ 3 answers. $\eta$ with 95\% bootstrap interval;
$E$ = range-normalised effect. Only factors whose interval excludes zero are listed.}
\label{tab:diag}
\begin{tabular}{lcccc}
\toprule
& \multicolumn{2}{c}{Qwen3-0.6B} & \multicolumn{2}{c}{Qwen3-1.7B} \\
\cmidrule(lr){2-3}\cmidrule(lr){4-5}
Factor & $\eta$ [95\% CI] & $E$ & $\eta$ [95\% CI] & $E$ \\
\midrule
steps          & 0.60 [0.49, 0.71] & 2.99 & 0.37 [0.26, 0.48] & 1.87 \\
n\_div         & 1.20 [0.97, 1.47] & 1.20 & 1.10 [0.89, 1.34] & 1.10 \\
log10\_max     & 0.29 [0.13, 0.45] & 0.76 & \multicolumn{2}{c}{not significant} \\
n\_mul         & 0.36 [0.15, 0.60] & 0.72 & 0.27 [0.08, 0.48] & 0.53 \\
n\_distractors & 0.34 [0.18, 0.52] & 0.34 & 0.33 [0.15, 0.50] & 0.33 \\
\bottomrule
\end{tabular}
\end{table}

\paragraph{Match quality.} For 0.6B, all 9{,}977 requested targeted problems were matched (mean $|\Delta
\text{logit}| = 0.006$, KS $= 0.034$). For 1.7B the off-target pool was thinner in the difficulty band, and the
match was looser (mean $|\Delta\text{logit}| = 0.32$, KS $= 0.27$); we return to this in
Section~\ref{sec:limits}.

\subsection{Targeting versus difficulty}
Table~\ref{tab:main} gives accuracy before and after training. Table~\ref{tab:primary} gives the primary
comparison, fixed before the runs: targeted minus matched control at 3:1.

\begin{table}[t]
\centering
\caption{Accuracy (mean of 3 seeds). In brackets: change versus the untrained model on the same problems. Every
GSM8K and GSM-Symbolic change has its 95\% interval below zero. ``Probes'' are generated problems from the training
families (train) or the two held-out families (held-out); ``$>4$'' is the held-out slice with more than four steps.}
\label{tab:main}
\small
\begin{tabular}{llccccc}
\toprule
Model & Arm & GSM8K & GSM-Sym. & Probes train & Probes held-out & Held-out $>4$ \\
\midrule
\multirow{6}{*}{0.6B}
& untrained          & 0.638 & 0.570 & 0.384 & 0.324 & 0.18 \\
& real\_only         & 0.458 ($-$0.20) & 0.415 ($-$0.16) & 0.237 ($-$0.13) & 0.247 ($-$0.10) & 0.09 \\
& untargeted 3:1     & 0.447 ($-$0.21) & 0.418 ($-$0.16) & 0.941 (+0.58) & 0.760 (+0.41) & 0.71 \\
& matched ctrl 3:1   & 0.448 ($-$0.21) & 0.408 ($-$0.17) & 0.881 (+0.52) & 0.760 (+0.41) & 0.65 \\
& targeted 3:1       & 0.436 ($-$0.22) & 0.407 ($-$0.17) & 0.795 (+0.43) & 0.602 (+0.25) & 0.54 \\
& targeted 9:1       & 0.422 ($-$0.24) & 0.382 ($-$0.19) & 0.869 (+0.51) & 0.704 (+0.36) & 0.66 \\
\midrule
\multirow{6}{*}{1.7B}
& untrained          & 0.815 & 0.773 & 0.625 & 0.583 & 0.41 \\
& real\_only         & 0.666 ($-$0.15) & 0.625 ($-$0.14) & 0.514 ($-$0.13) & 0.499 ($-$0.07) & 0.34 \\
& untargeted 3:1     & 0.660 ($-$0.15) & 0.600 ($-$0.16) & 0.973 (+0.33) & 0.926 (+0.36) & 0.89 \\
& matched ctrl 3:1   & 0.652 ($-$0.16) & 0.578 ($-$0.18) & 0.946 (+0.30) & 0.919 (+0.35) & 0.87 \\
& targeted 3:1       & 0.663 ($-$0.15) & 0.593 ($-$0.17) & 0.963 (+0.32) & 0.872 (+0.30) & 0.86 \\
& targeted 9:1       & 0.656 ($-$0.16) & 0.568 ($-$0.19) & 0.963 (+0.32) & 0.894 (+0.32) & 0.87 \\
\bottomrule
\end{tabular}
\end{table}

\begin{table}[t]
\centering
\caption{Targeted minus matched control at 3:1: paired difference in accuracy with 95\% bootstrap interval and
$p$-value.}
\label{tab:primary}
\begin{tabular}{lcc}
\toprule
Benchmark & Qwen3-0.6B & Qwen3-1.7B \\
\midrule
GSM8K             & $-0.012$ [$-0.039$, $+0.014$], $p=0.39$ & $+0.011$ [$-0.017$, $+0.039$], $p=0.42$ \\
GSM-Symbolic      & $-0.002$ [$-0.037$, $+0.032$], $p=0.93$ & $+0.015$ [$-0.020$, $+0.050$], $p=0.41$ \\
Probes, train     & $-0.086$ [$-0.143$, $-0.030$], $p=0.003$ & $+0.017$ [$-0.010$, $+0.049$], $p=0.24$ \\
Probes, held-out  & $-0.158$ [$-0.225$, $-0.089$], $p<0.001$ & $-0.047$ [$-0.094$, $-0.000$], $p=0.04$ \\
\bottomrule
\end{tabular}
\end{table}

\paragraph{No benefit on real problems.} On GSM8K and GSM-Symbolic the targeted and matched-control arms are
indistinguishable for both models; so are targeted and untargeted (0.6B: $-0.011$ and $-0.012$; 1.7B: $+0.003$
and $-0.007$; all intervals include zero).

\paragraph{The control wins on held-out generated problems.} On problems from the two held-out families the
targeted arm is below the control for both models, and for 0.6B it is below the control even on the slice it was
aimed at (problems with more than four steps: 0.54 versus 0.65). More synthetic data helps the targeted arm on
probes (9:1 above 3:1 on every probe measure) but does not close the gap to untargeted data.

\paragraph{Targeting does reduce the targeted effect.} Refitting the LLTM after training shows that the targeted arm
lowered the cost of an extra step far more than the other arms: for 0.6B, $\Delta\eta_{\text{steps}} = -0.98$
(targeted 3:1) and $-0.72$ (9:1), against $+0.05$ (untargeted) and $+0.23$ (matched control). The other arms instead
reduced the costs of multiplication and division ($\Delta\eta$ between $-0.90$ and $-1.76$, against $-0.01$ and
$-0.14$ for targeted 3:1). Long problems also contain multiplications and divisions, so the broader training data
produced higher accuracy on them even though it made ``more steps'' no less costly. For 1.7B the refit is poorly
identified, because the trained models answer 87--97\% of probes correctly, and we do not interpret it.

\subsection{Why every fine-tuned arm loses accuracy on GSM8K}
\label{sec:cost}
All arms lose 14--24 points on GSM8K and GSM-Symbolic, including real\_only, whose training data contains no
generated problems at all. Before the main runs, a development study on 0.6B (200 held-out GSM8K training
problems, one seed, exploratory) showed that the loss does not depend on the training settings we varied: a
four-times lower learning rate, a single epoch, or a larger share of real problems all left it between $-0.20$ and
$-0.24$. Rewriting every training solution into a longer fixed template made it worse ($-0.28$), because the model
began to repeat the template until it hit the token limit. Only replacing the training text with the base model's
own verified solutions reduced it, to $-0.115$ when 65\% of the tokens were model-written and $-0.08$ at 85\%, but
this also removed most of the gain on generated probes ($+0.26$ fell to $+0.10$ and $+0.08$).

To locate the loss, we compared each fine-tuned model's answers with the untrained model's answers to the same 300
GSM8K problems (Table~\ref{tab:cost}). Answer extraction is not the cause: almost no answer of either model lacks an
extractable final number. The fine-tuned models write much shorter answers, about four lines instead of 25--30, in
the format of the training solutions. The untrained model typically begins by listing every given quantity and then
solves one step per heading; the fine-tuned models skip that list. On problems that changed from correct to wrong,
the largest share of new errors are planning errors (a quantity left out, or a relation such as ``ten times more
than'' misread), followed by arithmetic slips in short inline sums. For 0.6B the rate of false equations rises
(from about 6 to 6--14 per 100 equations); for 1.7B it barely changes, so its loss is almost entirely in planning.

\begin{table}[t]
\centering
\caption{Paired answer-level comparison on the same 300 GSM8K test problems: untrained model versus each of the 15
fine-tuned models (5 arm settings $\times$ 3 seeds). ``Lost'' = answers correct before and wrong after training.
Slip rates are approximate, because equations written in \LaTeX{} by the untrained model are parsed less reliably.}
\label{tab:cost}
\small
\begin{tabular}{lcccc}
\toprule
& \multicolumn{2}{c}{Qwen3-0.6B} & \multicolumn{2}{c}{Qwen3-1.7B} \\
\cmidrule(lr){2-3}\cmidrule(lr){4-5}
& untrained & fine-tuned & untrained & fine-tuned \\
\midrule
Accuracy                      & 0.660 & 0.39--0.49 & 0.813 & 0.65--0.70 \\
Median answer length (chars)  & 728   & $\approx$215 & 808 & $\approx$210 \\
Median non-empty lines        & 25    & 4 & 30 & 4 \\
False equations per 100       & 5.8   & 6--14 & 5.0 & 4--9 \\
Lost / gained answers (pooled) & \multicolumn{2}{c}{1{,}222 / 241} & \multicolumn{2}{c}{915 / 222} \\
Lost answers: plan error      & \multicolumn{2}{c}{50\%} & \multicolumn{2}{c}{55\%} \\
Lost answers: arithmetic slip & \multicolumn{2}{c}{32\%} & \multicolumn{2}{c}{24\%} \\
Lost answers: distractor use  & \multicolumn{2}{c}{9\%}  & \multicolumn{2}{c}{12\%} \\
Lost answers: unverifiable    & \multicolumn{2}{c}{9\%}  & \multicolumn{2}{c}{9\%} \\
\bottomrule
\end{tabular}
\end{table}

Our reading is that supervised fine-tuning on short reference solutions mainly teaches the output format
\citep{zhou2023lima}, and in doing so overrides the longer reasoning format the model acquired in its own
post-training. That format, in particular the habit of restating the given facts first, is part of what made the
untrained model accurate. This is consistent with reports that self-generated text counters forgetting by matching
the model's style \citep{sun2026talking}. Because all arms pay the same cost, the comparison between arms remains
fair, but the cost limits the practical value of this training recipe.

% ---------------------------------------------------------------------------------------------------------------
\section{Discussion}

\paragraph{What the null result means.} In our setting the benefit of weakness-targeted data on real benchmarks is
indistinguishable from the benefit of equally hard data that does not target the weakness, and from untargeted
generated data. Studies that compare targeted data only with random data cannot tell these apart. The narrow
diagnostic measure (the targeted factor's LLTM weight) moved as intended, which shows that the intervention worked
on its own terms; it simply did not matter more than the broader skills the other arms improved.

\paragraph{Why aiming at one factor can hurt.} Real and generated long problems rarely fail for one reason. By
concentrating on the number of steps, the targeted arm saw fewer multiplications and divisions than the control,
and these remained costly. A single-factor target therefore trades one weakness for others. A natural next step is
to spread the budget over all significant factors in proportion to their effect; we are running this extension and
will report it separately.

\paragraph{Practical lesson.} On consumer hardware, a 0.6B or 1.7B model can be diagnosed, fine-tuned and re-evaluated
with five arm settings and three seeds in one to two days, but supervised training on short solution text costs far more general accuracy than targeting can
win back. Methods that keep the model's own reasoning style, such as rejection-sampled self-solutions with more
attempts on hard problems \citep{tong2024dart} or on-policy reinforcement learning with the generator as an exact
reward \citep{shenfeld2025razor,liang2025sws}, are the more promising direction.

\section{Limitations}
\label{sec:limits}
\begin{itemize}
  \item \textbf{Match quality for 1.7B.} The 1.7B control was matched less closely (KS 0.27) than the 0.6B control
  (KS 0.034). The 1.7B comparison should therefore be read as weaker evidence. \todo{decide whether to rerun the 1.7B
  single-weakness control with the enlarged pool used in later runs}
  \item \textbf{Ceiling effects.} The fine-tuned 1.7B models answer 87--97\% of generated probes correctly, which
  limits how well the arms can be separated on probes and makes the post-training LLTM refit unreliable.
  \item \textbf{Evaluation size.} 300 GSM8K and 200 GSM-Symbolic problems per model give paired intervals of about
  $\pm3$ points; an effect of one or two points would not be detected.
  \item \textbf{Error labels.} The automatic error classifier has not yet been validated against human labels
  \todo{report $\kappa$}, and GSM-Symbolic error types are not reported.
  \item \textbf{Scope.} One task family (arithmetic word problems), one model family (Qwen3), one training recipe
  (QLoRA supervised fine-tuning on short solutions), three seeds.
  \item \textbf{Synthetic text.} Generated problems read less naturally than human-written ones, and the probes share
  the generator with the training data, so part of the probe gains may be format learning; the held-out families
  reduce but do not remove this.
\end{itemize}

\section{Conclusion}
We built a controlled test of whether weakness-targeted synthetic data helps because it is targeted or because it
is hard. With exact skill profiles, a fitted LLTM and a one-to-one difficulty-matched control, two small Qwen3
models show the same picture: targeting reduces the targeted factor's cost but brings no gain on GSM8K or
GSM-Symbolic over equally hard off-target data, and it does worse on held-out generated problems. Separately, every
supervised arm loses substantial real-benchmark accuracy because it imitates terse solution text. Both results argue
that claims for targeted synthesis should be tested against a difficulty-matched control and alongside a check of
general accuracy.

\paragraph{Reproducibility.} Code, configurations, run records and per-answer outputs are in the project
repository \todo{public URL}. Paper runs: \texttt{80b37628b193} (0.6B) and \texttt{409305f3e2f2} (1.7B).

\bibliographystyle{plainnat}
\bibliography{references}

\end{document}
